% TOP_PROBLEM: {"id": 3351, "title": "Diophantine quintuple conjecture", "category_id": 1, "category": {"id": 1, "name": "number_theory", "display_name": "Number Theory", "description": "Properties of integers, prime numbers, Diophantine equations.", "slug": "number-theory", "order_index": 1, "created_at": "2026-07-31T15:26:25.670Z"}, "status": "open"}
% FIRST_POSTED: null
% ENTRY_KIND: statement_only
% Imported from: batch03.tex; UnsolvedMath ID 3351
\documentclass[11pt]{article}
\usepackage[margin=25mm]{geometry}
\usepackage{fontspec}
\setmainfont{Latin Modern Roman}
\usepackage{amsmath,amssymb,mathtools}
\usepackage{unicode-math}
\setmathfont{Latin Modern Math}
\usepackage{xurl,hyperref}
\hypersetup{colorlinks=true,urlcolor=blue}
\setcounter{secnumdepth}{0}
\setlength{\parindent}{0pt}
\setlength{\parskip}{5pt}
\setlength{\emergencystretch}{3em}
\newcommand{\sref}[2]{\hyperlink{source:#1:#2}{[S#2]}}
\newcommand{\cataloguescope}[1]{\paragraph{Recorded target.}#1}
\begin{document}
% UnsolvedMath ID 3351; source catalogue code OPG-16555
\setcounter{section}{3350}
\section{Diophantine quintuple conjecture}\label{3351}
{\small\sffamily UnsolvedMath ID 3351 · Number Theory}
\subsection{Definitions and mathematical statement}

\paragraph{Shared notation.}$\mathbb N=\{1,2,\ldots\}$, and $\mathbb Z,\mathbb Q,\mathbb R,\mathbb C$ denote the integers, rationals, reals and complex numbers. Any other conventions are specified locally.


A Diophantine $m$-tuple is a set of $m$ distinct positive integers such that the product of every two distinct members, plus one, is a perfect square. For a Diophantine triple $\{a,b,c\}$, let $r,s,t$ be the positive integers with $r^2=ab+1$, $s^2=ac+1$, $t^2=bc+1$, and define its regular extension\[ d_+=a+b+c+2abc+2rst. \] This is the formula in Dujella's original survey; the shorter $2bc$ term printed on the Garden page is a transcription error.
\paragraph{Statement recorded in the supplied source.}
(1) Is there no Diophantine quintuple of positive integers? (2) If $\{a,b,c,d\}$ is a Diophantine quadruple with $d>\max\{a,b,c\}$, must $d=d_+$? The second assertion is the stronger regular-extension conjecture from Dujella's survey. \sref{3351}{2} \sref{3351}{3}
\subsection{Short English statement}
The source asks both whether a Diophantine quintuple is impossible and whether every positive-integer quadruple is the regular extension of its three smaller members.
\subsection{Sources}
\begin{description}
\item[\hypertarget{source:3351:1}{S1}] ULAM.AI and the UnsolvedMath Contributors, UnsolvedMath: A Curated Collection of Open Mathematics Problems, record 3351 (OPG-16555). CC BY 4.0. \url{https://huggingface.co/datasets/ulamai/UnsolvedMath}.
\item[\hypertarget{source:3351:2}{S2}] Open Problem Garden, Diophantine quintuple conjecture, node 16555, original problem page. Original question and full discussion. \url{https://www.openproblemgarden.org/op/diophantine_quintuple_does_not_exist}.
\item[\hypertarget{source:3351:3}{S3}] Andrej Dujella, Diophantine quintuple conjecture, Section 2 of the author survey Diophantine m-tuples, displayed d+ formula and Conjecture 2.1, original author website source. \url{https://raw.githubusercontent.com/dujella/dujella.github.io/master/quint.html}.
\end{description}
\label{end:3351}
\end{document}
